\subsection{调和综合}

设 A 是一个交换 Banach 代数。回顾：若 M 是 \(\mathsf{X}(\mathsf{A})\) 的一个子集，我们用 \(\Upsilon(\mathsf{M})\) 表示 M 中元素的核的交（参见 I, p.~30）；它是 A 的一个理想。

命题 4. --- 设 A 是一个无根的正则交换 Banach 代数。设 F 是 X(A) 的一个闭子集。使 V(I) = F 的 A 的理想 I 所组成的集合按包含关系排序后有一个最大元，即 Y(F)，还有一个最小元，即由 G(x) 具有与 F 不相交的紧支集的 x ∈ A 组成的集合 J。

按构造，\(\Upsilon(F)\) 是 A 的一个理想，满足 \(V(\Upsilon(F))\) 包含 F，并且它是具有这一性质的 A 的最大理想。由于 F 是闭的，存在 A 的理想 I 使 \(V(I) = F\)；于是有 \(I \subset \Upsilon(F)\)，从而 \(V(\Upsilon(F)) \subset V(I) = F\)，即 \(V(\Upsilon(F)) = F\)。这就证明了第一个断言。

集合 J 是 A 的理想且 V(J) 包含 F。我们来证明 V(J) = F。设 \(\chi \in \mathsf{X}(\mathsf{A})\) 不属于 F。取 \(\chi\) 的一个与 F 不相交的紧邻域 U (TG, I, p.~65, 命题 9 的推论)。依 I, p.~88 命题 1 的断言 (ii)，存在 \(x \in A\) 使 \(\mathcal{G}(x)\) 在 \(\chi\) 处等于 1 而在 U 之外等于 0。于是有 \(x \in J\)，因此 \(\chi \notin \mathrm{V}(\mathrm{J})\)。这表明 V(J) ⊂ F，从而 V(J) = F。

最后，设 I 是 A 的理想且 \(V(I) = F\)。我们证明 \(J \subset I\)。设 \(x \in J\)，并设 C 是 \(\mathcal{G}(x)\) 的支集；集合 C 是 \(X(A)\) 的一个与 F 不相交的紧子集。依命题 3，存在元素 \(u \in I\) 使 \(\mathcal{G}(u) = 1\) 在 C 上成立。于是有 \(\mathcal{G}(x) = \mathcal{G}(ux)\)，又因 A 无根，故 x = ux (I, p.~38, 命题 8)。从而有 \(x \in I\)，这表明 \(J \subset I\)。

推论 1. --- 设 A 是一个无根的正则交换 Banach 代数。设 J 是由 \(x \in \mathrm{A}\) 中使 \(\mathcal{G}(x)\) 具有紧支集的那些元素组成的集合。假设 \(\overline{\mathrm{J}} = \mathrm{A}\)。那么 A 的每个异于 A 的闭理想都含于一个正则极大理想。

若 I 是 A 的闭理想且不含于任何正则极大理想，则 \(V(I) = \varnothing\)，于是 \(I \supset J\)（命题 4 应用于 \(F = \varnothing\)），从而 \(I \supset \overline{J} = A\)。

推论 2. --- 设 A 是一个无根的正则交换 Banach 代数。设 \(x, y \in A\)。若 \(\mathcal{G}(x)\) 的支集是紧的，且含于由特征标 \(\chi\) 组成的集合之中且在其上 \(\mathcal{G}(y)(\chi) \neq 0\)，则 \(x\) 是 \(y\) 在 \(A\) 中的倍元。

设 I 是 A 的理想 Ay。于是 V(I) 是 \(\mathcal{G}(y)\) 的零点集。由于 \(\mathcal{G}(x)\) 的支集 F 是紧的且与 V(I) 不相交，我们有 \(x \in I\)（命题 4 应用于 F）。

定义 2. --- 设 A 是一个交换 Banach 代数。

设 I 是 A 的理想，\(x \in A\)，\(\chi \in \mathsf{X}'(\mathsf{A})\)。我们称 x 在 \(\chi\) 的邻域内属于 I，如果存在元素 \(y \in I\) 使 \(\mathcal{G}'(y)\) 与 \(\mathcal{G}'(x)\) 在 \(\chi\) 的一个邻域中重合。

若对每个 \(\chi \in \mathsf{X}'(\mathrm{A})\) 以及每个 \(x \in \mathrm{A}\)（\(\mathcal{G}'(x)\) 在 \(\chi\) 处为零），存在 A 中的序列 \((x_n)\) 使 \(x = \lim_{n \to \infty} x_n x\)，且每个 \(\mathcal{G}'(x_n)\) 在 \(\mathrm{V}_n\)（\(\chi\) 的一个邻域）中为零，则称 A 满足 Ditkin 条件。

注记. --- 设 A 是一个交换 Banach 代数。

\begin{enumerate}
\def\labelenumi{\arabic{enumi})}
\item
  若 \(\chi\) 使 \(\mathcal{G}'(x)\) 在 \(\chi\) 的一个邻域中为零，则 \(x\) 在 \(\chi\) 的邻域内属于 I。
\item
  若 x 在 \(\chi\) 的邻域内属于 I，且 \(y \in A\) 是任意元素，则 xy 在 \(\chi\) 的邻域内属于 I。
\item
  由特征标 \(\chi\) 组成的、使 \(x\) 在 \(\chi\) 的邻域内属于 I 的集合在 \(\mathrm{X}'(\mathrm{A})\) 中是开的。
\item
  假设 A 是正则且无根的。若 x 在每个特征标 \(\chi\) 的邻域内属于 I，其中 \(\chi \in \mathsf{X}'(\mathrm{A})\)，则 x 属于 I（I, p.~91 的推论 2 应用于函数 \(f = \mathcal{G}'(x)\)，以及 I, p.~38 的命题 8）。
\item
  假设 A 是正则的。设 I 是 A 的理想，并设 \(\chi\) 是 \(\mathsf{X}(\mathsf{A})\) 中使 \(\chi\notin\mathsf{V}(\mathsf{I})\) 的元素。那么 A 的每个元素 \(x\) 都在 \(\chi\) 的邻域内属于 I。事实上，依 I, p.~89 的定义 1，存在 \(z\in\mathsf{A}\) 使 \(\mathcal{G}'(z)\) 在 \(\chi\) 的一个邻域中等于 1，而在 V(I) 的一个邻域中等于 0。\(\mathcal{G}(z)\) 的支集是紧的，于是有 \(z\in\mathsf{I}\)（命题 4 应用于 V(I)）；因此 \(xz\in\mathsf{I}\)，且 \(\mathcal{G}'(xz)=\mathcal{G}'(x)\) 在 \(\chi\) 的一个邻域中成立。
\end{enumerate}

回顾：拓扑空间 X 的子空间 K 称为完全的，如果它是闭的且没有孤立点 (TG, I, p.~8)。

引理 1. --- 设 A 是一个无根的满足 Ditkin 条件的正则交换 Banach 代数。设 I 是 A 的一个闭理想，并设 x 是 \(\Upsilon(\mathrm{V}(\mathrm{I}))\) 的一个元素。设 K 是由 \(\chi \in \mathsf{X}'(\mathrm{A})\) 中使 x 不在 \(\chi\) 的邻域内属于 I 的特征标组成的集合。那么集合 K 是 \(\mathsf{X}'(\mathrm{A})\) 的一个完全子集。

用 G 表示 K 在 \(X'(A)\) 中的余集。集合 G 在 \(X'(A)\) 中是开的（注 3），故 K 是闭的。

用反证法，设 K 有一个孤立点 \(\chi_{0}\)。取 \(\chi_{0}\) 的一个使 \(U - \{\chi_{0}\}\subset G\) 的邻域 U。由于 x 不在 \(\chi_{0}\) 的邻域内属于 I，注 5 表明 \(\chi_{0}\in\mathrm{V(I)}\)。特别地，我们有 \(\chi_{0}(x)=0\)，因为 \(x\in\Upsilon(\mathrm{V(I)})\)。

我们将证明：存在 A 的元素 \(y\)，它在 \(\mathsf{X}'(\mathrm{A}) - \{\chi_0\}\) 的每一点的邻域内属于 I，不在 \(\chi_0\) 的邻域内属于 I，并且 \(\chi_0(y) = 0\)。

首先假设这样的元素 \(y\) 的存在性已经证明。由于 A 满足 Ditkin 条件，于是存在 A 中的序列 \((x_n)\) 使 \(x_n y\) 趋于 \(y\)，且每个 \(\mathcal{G}'(x_n)\) 在 \(\chi_0\) 的一个邻域中为零。于是对每个 \(n\)，元素 \(x_n y\) 在 \(X'(A)\) 的每一点的邻域内属于 I（注 1 与注 2），因此 \(x_{n}y \in I\)（注 4）。由于 I 是闭的，我们推出 \(y \in I\)，这与 y 不在 \(\chi_{0}\) 的邻域内属于 I 相矛盾。

剩下的就是证明 \(y\) 的存在性。若 \(\chi_0 \neq 0\)，依 I, p.~88 命题 1 的断言 (iii)，存在 \(u \in A\) 使 \(\mathcal{G}'(u)\) 在 \(\chi_0\) 的一个邻域中等于 1，而在 \(\mathrm{X}'(\mathrm{A}) - \mathrm{U}\) 的一个邻域中等于 0。设 \(y = ux\)。由于 \(x\) 对每个 \(\chi\)（这里 \(\chi \in \mathrm{U} - \{\chi_0\}\)）都在其邻域内属于 I，故对 \(y\) 亦然。此外，若 \(\chi \in \mathrm{X}'(\mathrm{A}) - \mathrm{U}\)，则 \(\mathcal{G}'(y)\) 在 \(\chi\) 的一个邻域中为零。于是（注 5）元素 \(y = ux\) 在每个 \(\chi \neq \chi_0\) 的邻域内属于 I。由于 \(\mathcal{G}'(y)\) 与 \(\mathcal{G}'(x)\) 在 \(\chi_0\) 的一个邻域中重合，\(\chi_0\) 属于 K 这一事实蕴含 \(y\) 不在 \(\chi_0\) 的邻域内属于 I。最后，我们有 \(\chi_0(y) = \chi_0(u)\chi_0(x) = 0\)。

若 \(\chi_0 = 0\)，类似地存在元素 \(v \in \mathrm{A}\) 使 \(\mathcal{G}'(v)\) 在 \(\chi_0\) 的一个邻域中为零，而在 \(\mathsf{X}'(\mathsf{A}) - \mathsf{U}\) 的一个邻域中等于 1；与前同理，我们由此推出元素 \(y = x - vx\) 在每个 \(\chi \neq \chi_0\) 的邻域内属于 I，它不在 \(\chi_0\) 的邻域内属于 I，且 \(\chi_0(y) = 0\)。

引理 2. — 设 X 是一个拓扑空间。设 F 与 D 是 X 的两个不相交的子空间，使得 F 是闭的而 D 是离散的。若 F 不含非空完全子空间，则 F ∪ D 亦然。

事实上，设 K 是 F ∪ D 的一个非空完全子空间。设 x 是 K 的一点。若 x 属于 D，则它在 D 中是孤立的，又因 F 是闭的，它在 F ∪ D 中也是孤立的。于是 x 在 K 中是孤立的，这与假设矛盾。

命题 5. — 设 A 是一个无根的满足 Ditkin 条件的正则交换 Banach 代数。设 I 是 A 的一个闭理想，使得 V(I) 的边界 F 不含非空完全集。那么 I = Y(V(I))，也就是说，I 是由使 G(x) 在 V(I) 上为零的 x ∈ A 组成的集合。特别地，若 V(I) 化为单点 χ，则 I = Ker(χ)。

我们有 \(\mathrm{I} \subset \Upsilon(\mathrm{V}(\mathrm{I}))\)。现设 \(x \in \Upsilon(\mathrm{V}(\mathrm{I}))\)。设 \(\mathrm{G}\) 是由特征标 \(\chi \in \mathsf{X}'(\mathbf{A})\) 组成的、使 \(x\) 属于 \(\mathrm{I}\)（在 \(\chi\) 的一个邻域内）的集合。它是开的，其余集 \(\mathrm{K}\) 是完全集（引理 1）。由于 \(\mathcal{G}'(x)\) 在 \(\mathrm{V}(\mathrm{I})\) 上为零，集合 \(\mathrm{G}\) 包含 \(\mathrm{V}(\mathrm{I}) \cup \{0\}\) 的内部（注 1）。依注 5，它还包含 \(\mathsf{X}(\mathbf{A}) - \mathsf{V}(\mathbf{I})\)。因此 \(\mathrm{K} = \mathsf{X}'(\mathbf{A}) - \mathsf{G}\) 含于 \(\mathrm{F}_0\)（即 \(\mathrm{V}(\mathrm{I}) \cup \{0\}\) 的边界）之中。

我们有 \(F_{0} \subset F \cup \{0\}\)。因此假设蕴含 \(F_{0}\) 不含非空完全集（引理 2）。于是由引理 1 可知完全集 K 是空集。从而 x 在每个 \(\chi \in \mathsf{X}'(\mathsf{A})\) 的邻域内属于 I，这意味着 \(x \in I\)（注 4）。这样就有 \(\Upsilon(\mathrm{V}(\mathrm{I})) \subset \mathrm{I}\)，证明完毕。

